\documentclass[12pt,a4paper]{article}
\usepackage[T1]{fontenc}
\usepackage[utf8]{inputenc}
\usepackage[brazil]{babel}
\usepackage{geometry}
\usepackage{xcolor}
\usepackage{fancyhdr}
\usepackage{titlesec}
\usepackage{parskip}
\usepackage{booktabs}
\usepackage{mdframed}
\usepackage{hyperref}
\usepackage{setspace}

\geometry{top=2.5cm,bottom=2.5cm,left=3cm,right=2.5cm}

\definecolor{ufamblue}{RGB}{0,51,102}
\definecolor{lightblue}{RGB}{230,240,255}
\definecolor{metabg}{RGB}{245,245,245}

\pagestyle{fancy}
\fancyhf{}
\rhead{\small\color{ufamblue} PPGEE/UFAM --- Seleção Mestrado 2026}
\lhead{\small\color{ufamblue} Redação Técnica Objetiva}
\cfoot{\thepage}
\renewcommand{\headrulewidth}{0.4pt}

%==============================================================================
\begin{document}
%==============================================================================

%-- Cabeçalho institucional ---------------------------------------------------
\begin{center}
  \textcolor{ufamblue}{\Large\bfseries Universidade Federal do Amazonas}\\[3pt]
  \textcolor{ufamblue}{\large Programa de Pós-Graduação em Engenharia Elétrica}\\[2pt]
  {\normalsize Edital N.\textordmasculine\ 023/2026-PROPESP/UFAM}\\[10pt]
  \rule{\linewidth}{1.5pt}\\[6pt]
  {\LARGE\bfseries Redação Técnica Objetiva}\\[4pt]
  {\large\bfseries Seleção Mestrado Acadêmico 2026}\\[6pt]
  \rule{\linewidth}{1.5pt}
\end{center}

\vspace{6mm}

%-- Dados do candidato --------------------------------------------------------
\begin{mdframed}[backgroundcolor=metabg,linecolor=ufamblue,linewidth=1pt,
                 innertopmargin=8pt,innerbottommargin=8pt]
\begin{tabular}{ll}
  \textbf{Candidato:}       & Pedro Victor dos Santos Oliveira \\[3pt]
  \textbf{Linha de pesquisa:} & Linha 1 --- Sistemas Inteligentes e Microeletrônica \\[3pt]
  \textbf{Tema escolhido:}  & Aprendizado Profundo (\textit{deep learning}) \\[3pt]
  \textbf{Data:}            & \today \\
\end{tabular}
\end{mdframed}

\vspace{8mm}

%-- Redação -------------------------------------------------------------------
\begin{mdframed}[backgroundcolor=lightblue,linecolor=ufamblue,linewidth=1.2pt,
                 innertopmargin=14pt,innerbottommargin=14pt,
                 innerleftmargin=14pt,innerrightmargin=14pt]

{\setstretch{1.5}

O avanço do aprendizado profundo elevou substancialmente o desempenho em
tarefas de visão computacional e processamento digital de imagens. Contudo, a
execução desses modelos ainda depende, na maioria dos casos, de infraestrutura
computacional de alto custo, tornando sua adoção inviável em dispositivos
móveis e sistemas embarcados com recursos limitados.

A inteligência artificial embarcada, ou \textit{Edge AI}, propõe deslocar o
processamento dos servidores para a borda da rede, reduzindo latência, consumo
de largura de banda e dependência de conectividade. Técnicas como quantização
de pesos, poda de redes neurais e destilação de conhecimento permitem
comprimir modelos de alta capacidade sem perda expressiva de acurácia,
viabilizando sua execução em microcontroladores e \textit{smartphones} de
entrada.

Em regiões com infraestrutura de comunicação precária, como parte da Amazônia,
essa abordagem é estratégica: sistemas de diagnóstico por imagem, monitoramento
ambiental e reconhecimento de padrões agrícolas podem operar de forma autônoma,
sem acesso contínuo à nuvem. A convergência entre PDI e \textit{Edge AI}
representa uma fronteira de pesquisa com alto impacto social e científico.

O ingresso no Mestrado do PPGEE/UFAM representa a oportunidade de investigar
modelos eficientes de visão computacional voltados ao processamento embarcado,
contribuindo com soluções tecnológicas adaptadas à realidade regional amazônica.

}
\end{mdframed}

\vspace{8mm}

%-- Metadados da redação ------------------------------------------------------
\begin{center}
\begin{tabular}{lc}
\toprule
\textbf{Métrica} & \textbf{Valor} \\
\midrule
Caracteres com espaços & 1.413 \\
Limite mínimo exigido  & 500 \\
Limite máximo exigido  & 1.500 \\
Situação               & \textcolor{green!50!black}{\bfseries Dentro do limite} \\
\midrule
Palavras               & 195 \\
\bottomrule
\end{tabular}
\end{center}

\vspace{6mm}

%-- Critérios de avaliação ----------------------------------------------------
\begin{mdframed}[backgroundcolor=metabg,linecolor=gray!50,linewidth=0.8pt,
                 innertopmargin=8pt,innerbottommargin=8pt]
\textbf{\small Critérios de avaliação (Edital, item 4.1.3):}\\[4pt]
\small
\begin{tabular}{lc}
Adequação e domínio do tema         & 50\% \\
Organização textual e coerência     & 15\% \\
Objetividade e capacidade de síntese & 15\% \\
Estrutura (intro / arg. / discussão / conclusão) & 10\% \\
Grafia, acentuação e gramática      & 10\% \\
\end{tabular}
\end{mdframed}

\vspace{6mm}

%-- Nota de revisão -----------------------------------------------------------
\begin{mdframed}[backgroundcolor=yellow!15,linecolor=orange!60,linewidth=1pt,
                 innertopmargin=8pt,innerbottommargin=8pt]
\small\textbf{Nota para revisão:} reescreva a redação com suas próprias
experiências antes de submeter. Substitua trechos genéricos por referências
a projetos reais que você realizou (ex.: trabalhos de PDI, iniciação
científica, TCC). O formulário eletrônico do PPGEE aceita apenas texto
digitado diretamente no campo --- não é permitido anexar documento.
\end{mdframed}

\end{document}
